
$f$
étant le morphisme structural $X\longrightarrow S$.

Les $K^{\bullet}$ considérés seront ceux qui sont définis à l'aide des faisceaux localement libres de type fini. On notera $E$, $F$, $\xi$, les classes de $\underline{E}$, $\underline{F}$, $\underline{O}_X(1)$. On pose $x=1-\xi^{-1}$, $x'=\xi-1=\xi x$, et on note $\bar x$, $\bar x'$ les classes de $x$, $x'$ dans $\Gr^1K^{\bullet}(X)$. Comme $\xi-1$ et $x$ sont dans $\Fil^1K^{\bullet}(X)$, on a $(\xi-1)x\in\Fil^2K^{\bullet}(X)$, i.e. $x\equiv \xi x=x'$ mod. $\Fil^2K^{\bullet}(X)$ i.e. $\bar x'= \bar x$

5.2.
Si $\underline{G}$ est un $\underline{O}_X$-Module localement libre de type fini, les $R^if_*(\underline{G})$ ne sont pas en général localement libres de type fini sur $S$, si bien qu'on ne peut définir un homomorphisme $f_*:K^{\bullet}(X)\to K^{\bullet}(S)$ comme on le fait pour les $K^{\bullet}$ définis à l'aide des complexes parfaits. Le théorème 1.1 permet de tourner cette difficulté de la façon suivante : si $y\in K^{\bullet}(X)$, on peut, d'une seule manière, écrire :
\begin{equation}\tag{5.2.1}
  y=a_0+a_1\xi^{-1}+\ldots+a_r\xi^{-r},\qquad \text{avec } a_i\in K^{\bullet}(S);
\end{equation}
on pose alors :
\begin{equation}\tag{5.2.2}
  f_*(y)=a_0 .
\end{equation}

Il est clair que l'homomorphisme ainsi défini est additif, et vérifie la formule de projection, i.e. est $K^{\bullet}(S)$-linéaire.

On vérifie les relations suivantes :
\begin{equation}\tag{5.2.3}
  f_*(1)=1;\qquad
  f_*(\xi^n)=0\quad\text{pour }-r\le n<0;\qquad
  f_*(\xi^n)=\sigma_n,
  \quad 1\le n .
\end{equation}

Pour vérifier le dernier groupe de relations, on procède par récurrence sur $n$. Supposons que pour $p\le n-1$, $f_*(\xi^p)=\sigma_p$; la relation (1.9.1), multipliée par une puissance convenable de $\xi$, donne :
\[
  \sum_{i=0}^{r+1}(-1)^i\lambda^i(E)\xi^{\,n-i}=0 .
\]
On peut lui appliquer $f_*$ :
